%% kapa: exam=1 round=25 year=2014 period=2 form=A subject=accounting subject_ko=회계학 seq=39 number=79
%% source: Q-Net 감정평가사 기출문제 (cst00309) artlSeq=5101207 file=extracted/1008/5101207/2014년 제25회 감정평가사 1차시험 기출문제/2교시(A형).hwp
%% conversion: kordoc 4.24.51 -> tools/gen1.py (auto); review=proofread-sonnet; flags=none
\begin{question}{79}
(주)감평은 A제품과 B제품을 생산ㆍ판매하고 있으며, 다음 연도 예산손익계산서는 다음과 같다.
\par{\small\centering
\begin{tabular}{@{}l@{\hspace{2em}}r@{\hspace{2em}}r@{}}
\hline
 & \uline{A제품} & \uline{B제품} \\
매출액 & ￦4,000 & ￦2,000 \\
변동원가 & 1,500 & 1,200 \\
고정원가 & 2,000 & 1,400 \\
영업이익(손실) & ￦500 & (￦600) \\
판매량 & 2,000단위 & 2,000단위 \\
\hline
\end{tabular}\par}
회사는 영업손실을 초래하고 있는 B제품의 생산을 중단하고자 한다. B제품의 생산을 중단하면, A제품의 연간 판매량이 1,000단위만큼 증가하고 연간 고정원가 총액은 변하지 않는다. 이 경우 회사 전체의 영업이익은 얼마나 증가(혹은 감소)하는가? (단, 기초 및 기말 재고자산은 없다.) \par
\begin{choices}
  \item ￦175 감소
  \item ￦450 증가
  \item ￦650 감소
  \item ￦1,250 증가
  \item ￦1,425 증가
\end{choices}
\end{question}
\answer{2}
